% ===================================================================
%  Семинар 7. Содержательная часть — единый источник для обоих файлов.
%
%  Состав и порядок задач — по «Семинар 7.docx» («Автокорреляция»):
%      теория — метод рядов, тест Дарбина–Уотсона (таблица d_L, d_U);
%      задача 6.6 — DW и коррекция AR(1) (данные — «Семинар 7 (приложение).xlsx»);
%      задача 6.3 — обобщение Кохрейна–Оркатта на AR(2) («если разберём К–О»).
%  Источник: Сборник задач к начальному курсу эконометрики (РАНХиГС),
%  гл. 6, с. 176–177 и 179–181.
%
%  Задача 6.6 разбиралась и в семинаре 6; здесь она дополнена методом рядов.
%  В конспекте формула границ метода рядов записана как M(k) ± z·D(k);
%  верно M(k) ± z·sqrt(D(k)) — D(k) здесь дисперсия.
%
%  Раздел «Основные формулы» и всё внутри \begin{solution}...\end{solution}
%  попадает только в seminar_07_solutions.pdf.
% ===================================================================

\seminartitle{Семинар 7}{Автокорреляция ошибок}

% ===================================================================
\ifsolutions
% Editable vector illustrations for seminar 7 (TikZ).
% Coordinates of the residual plots (problem 6.6), computed from the task data.
\newcommand{\olsTimeCoordinates}{(1.000000,-0.070722) (2.000000,-0.411736) (3.000000,-0.419468) (4.000000,-0.106693) (5.000000,-0.637064) (6.000000,-0.136421) (7.000000,0.065031) (8.000000,-0.022058) (9.000000,0.143031) (10.000000,0.031247) (11.000000,0.121186) (12.000000,0.179866) (13.000000,-0.004311) (14.000000,0.066977) (15.000000,-0.134230) (16.000000,0.051990) (17.000000,0.056645) (18.000000,-0.082867) (19.000000,-0.061657) (20.000000,0.277906) (21.000000,0.099646) (22.000000,0.571734) (23.000000,0.139581) (24.000000,0.282384)}
\newcommand{\arTimeCoordinates}{(2.000000,-0.399915) (3.000000,-0.264760) (4.000000,0.066042) (5.000000,-0.612357) (6.000000,0.196608) (7.000000,0.140855) (8.000000,-0.066946) (9.000000,0.115915) (10.000000,-0.025101) (11.000000,0.077717) (12.000000,0.131712) (13.000000,-0.080380) (14.000000,0.081664) (15.000000,-0.203237) (16.000000,0.141121) (17.000000,0.012521) (18.000000,-0.093095) (19.000000,-0.027291) (20.000000,0.262160) (21.000000,0.007329) (22.000000,0.471971) (23.000000,-0.106096) (24.000000,0.173564)}

\newenvironment{seminarvisual}{\par\medskip\noindent\begin{minipage}{\linewidth}\centering}
  {\end{minipage}\par\medskip}
\newcounter{visualfig}
\newcommand{\visualcaption}[1]{\par\smallskip\refstepcounter{visualfig}%
  {\small\textbf{Рисунок~\thevisualfig.}\ #1\par}}

\newcommand{\dwvisual}{%
\begin{seminarvisual}
\begin{tikzpicture}[x=3cm,y=1cm,font=\small]
  \fill[accent!15] (0,0) rectangle (1.27,.6);
  \fill[notebar!25] (1.27,0) rectangle (1.45,.6);
  \fill[solbar!12] (1.45,0) rectangle (4,.6);
  \draw[->] (0,0)--(4.1,0) node[right] {DW};
  \draw[dashed] (1.27,0)--(1.27,.8);
  \draw[dashed] (1.45,0)--(1.45,.8);
  \node[anchor=north east] at (1.27,-.12) {$d_L=1{,}27$};
  \node[anchor=north west] at (1.45,-.12) {$d_U=1{,}45$};
  \node[below] at (0,0) {0}; \node[below] at (4,0) {4};
  \node at (.60,.3) {Отвергаем $H_0$};
  \node at (2.75,.3) {Не отвергаем $H_0$};
  \node[anchor=south west] at (1.5,1.25) {Зона неопределённости};
  \draw[->,notebar] (1.65,1.2)--(1.36,.65);
  \draw[->,very thick,accent] (1.103,1.08)--(1.103,.64);
  \node[anchor=south east,text=accent] at (1.13,1.08) {$\mathrm{DW}=1{,}103$};
\end{tikzpicture}
\visualcaption{Задача 6.6: $H_0\!:\rho=0$ против $H_1\!:\rho>0$,
уровень 5\,\%. Наблюдаемая статистика попадает в область отвержения.
Правая область относится только к этой односторонней проверке.}
\end{seminarvisual}}

\newcommand{\arresidualvisual}{%
\begin{seminarvisual}
\begin{tikzpicture}[x=.235cm,y=2.3cm,font=\small]
  \foreach \shift/\title in {0/{Исходные остатки $e_t$},30/{Инновации после AR(1)}} {
    \begin{scope}[xshift=\shift*.235cm]
      \node[font=\small\bfseries] at (12.5,.88) {\title};
      \draw[->] (0,-.65)--(0,.70);
      \draw[gray,dashed] (0,0)--(24,0);
      \draw[->] (0,-.65)--(25,-.65) node[right] {$t$};
      \foreach \xx in {1,6,12,18,24}
        \draw (\xx,-.65)--(\xx,-.69) node[below] {\xx};
      \foreach \yy/\lab in {-.5/{-0{,}5},0/0,.5/{0{,}5}}
        \node[left] at (0,\yy) {$\lab$};
      \ifnum\shift=0
        \draw[accent,thick] plot[mark=*,mark size=1.7pt] coordinates {\olsTimeCoordinates};
      \else
        \draw[solbar,thick] plot[mark=*,mark size=1.7pt] coordinates {\arTimeCoordinates};
      \fi
    \end{scope}
  }
\end{tikzpicture}
\visualcaption{Данные задачи 6.6, одинаковый масштаб по вертикали.
Слева $\mathrm{DW}=1{,}103$, справа $2{,}162$.
Справа показаны $\widehat\eta_t=\widehat\varepsilon_t-
\widehat\rho\widehat\varepsilon_{t-1}$, поэтому ряд начинается с $t=2$.
График дополняет диагностику, но не доказывает независимость.}
\end{seminarvisual}}

\newcommand{\intervalvisual}{%
\begin{seminarvisual}
\begin{tikzpicture}[x=14cm,y=1cm,font=\small]
  % Plot x = beta + 1.1 to keep node positions compact.
  \draw[->] (0,0)--(.63,0) node[right] {$\beta_2$};
  \foreach \xx/\lab in {0/{-1{,}1},.1/{-1{,}0},.2/{-0{,}9},.3/{-0{,}8},.4/{-0{,}7},.5/{-0{,}6},.6/{-0{,}5}}
    \draw (\xx,0)--(\xx,-.07) node[below] {$\lab$};
  \node[anchor=east,align=right] at (-.02,1.65) {МНК};
  \node[anchor=east,align=right] at (-.02,.65) {С AR(1)};
  \draw[accent,line width=1.3pt] (.086,1.65)--(.556,1.65);
  \draw[solbar,line width=1.3pt] (.075,.65)--(.462,.65);
  \foreach \xx in {.086,.556} \draw[accent] (\xx,1.55)--(\xx,1.75);
  \foreach \xx in {.075,.462} \draw[solbar] (\xx,.55)--(\xx,.75);
  \fill[accent] (.3209,1.65) circle (2.7pt);
  \fill[solbar] (.2688,.65) circle (2.7pt);
  \node[above] at (.086,1.75) {$-1{,}014$};
  \node[above] at (.556,1.75) {$-0{,}544$};
  \node[above] at (.075,.75) {$-1{,}025$};
  \node[above] at (.462,.75) {$-0{,}638$};
\end{tikzpicture}
\visualcaption{Точки — оценки $\beta_2$, отрезки — 95\,\%-ные интервалы.
Интервал МНК построен при нарушенной предпосылке независимости;
интервал с AR(1) — приближённый.}
\end{seminarvisual}}

\newcommand{\autocorrvisual}{%
\begin{seminarvisual}
\begin{tikzpicture}[x=.3cm,y=1.1cm,font=\small]
  \foreach \shift/\title in {0/{Нет автокорреляции},26/{Положительная автокорреляция}} {
    \begin{scope}[xshift=\shift*.3cm]
      \node[font=\small\bfseries] at (10.5,1.25) {\title};
      \draw[->] (0,-1.05)--(0,1.05) node[left] {$e_t$};
      \draw[->,gray] (0,0)--(21.5,0) node[right] {$t$};
      \ifnum\shift=0
        \draw[accent,thick] plot[mark=*,mark size=1.5pt] coordinates
          {(1,.5) (2,-.6) (3,.2) (4,.7) (5,-.3) (6,-.8) (7,.4) (8,-.2) (9,.6) (10,-.5)
           (11,.1) (12,.8) (13,-.7) (14,.3) (15,-.4) (16,.5) (17,-.1) (18,-.6) (19,.7) (20,-.3)};
      \else
        \draw[accent,thick] plot[mark=*,mark size=1.5pt] coordinates
          {(1,.3) (2,.5) (3,.7) (4,.6) (5,.4) (6,.1) (7,-.2) (8,-.5) (9,-.7) (10,-.8)
           (11,-.6) (12,-.3) (13,.1) (14,.4) (15,.6) (16,.8) (17,.5) (18,.2) (19,-.2) (20,-.4)};
      \fi
    \end{scope}
  }
\end{tikzpicture}
\visualcaption{Схематический пример. Слева знаки остатков меняются часто и
беспорядочно. Справа соседние остатки похожи: остатки долго остаются
одного знака, рядов одинаковых знаков мало, разности $e_t-e_{t-1}$ малы.}
\end{seminarvisual}}

\newcommand{\dwzonesvisual}{%
\begin{seminarvisual}
\begin{tikzpicture}[x=3.6cm,y=1cm,font=\small]
  \fill[accent!15]  (0,0)   rectangle (.9,.7);
  \fill[notebar!25] (.9,0) rectangle (1.4,.7);
  \fill[solbar!12]  (1.4,0) rectangle (2.6,.7);
  \fill[notebar!25] (2.6,0) rectangle (3.1,.7);
  \fill[accent!15]  (3.1,0) rectangle (4,.7);
  \draw[->] (0,0)--(4.12,0) node[right] {DW};
  \foreach \xx/\lab in {0/0,.9/d_L,1.4/d_U,2/2,2.6/4-d_U,3.1/4-d_L,4/4}
    \draw (\xx,0)--(\xx,-.1) node[below] {$\lab$};
  \node[align=center] at (.45,.35) {положительная\\[-2pt]автокорреляция};
  \node[align=center] at (2,.35) {автокорреляции нет};
  \node[align=center] at (3.55,.35) {отрицательная\\[-2pt]автокорреляция};
  \node[anchor=south,align=center] at (1.15,.8) {вывода\\[-2pt]нет};
  \node[anchor=south,align=center] at (2.85,.8) {вывода\\[-2pt]нет};
  \node[anchor=south,text=gray] at (.45,.8) {$\rho\approx 1$};
  \node[anchor=south,text=gray] at (2,.8) {$\rho\approx 0$};
  \node[anchor=south,text=gray] at (3.55,.8) {$\rho\approx -1$};
\end{tikzpicture}
\visualcaption{\label{fig:dwzones}Пять зон для статистики DW. Границы $d_L$ и $d_U$ берутся
из таблицы по $n$ и числу регрессоров без константы; правая половина
симметрична левой относительно $\mathrm{DW}=2$.}
\end{seminarvisual}}

\newcommand{\coloopvisual}{%
\begin{seminarvisual}
\begin{tikzpicture}[font=\small,box/.style={draw=accent,rounded corners=2pt,
  fill=accent!5,align=center,text width=3.05cm,minimum height=1.05cm,inner sep=2pt}]
  \node[box] (a) at (0,0) {1) МНК исходной\\модели, $e_t$};
  \node[box] (b) at (3.85,0) {2) оценка $\widehat\rho$\\по остаткам $e_t$};
  \node[box] (c) at (7.7,0) {3) квазиразности:\\МНК $\to\widehat\beta_1$, $\widehat\beta_2$};
  \node[box,text width=3.55cm] (d) at (11.8,0) {4) новые остатки\\$e_t=y_t-\widehat\beta_1-\widehat\beta_2x_t$};
  \draw[->,thick,accent] (a)--(b);
  \draw[->,thick,accent] (b)--(c);
  \draw[->,thick,accent] (c)--(d);
  \draw[->,thick,notebar] (d.south) -- ++(0,-.38) -| node[pos=.25,fill=white,inner sep=1.5pt]
    {пока $\widehat\rho$ меняется} (b.south);
\end{tikzpicture}
\visualcaption{Итерации Кохрейна--Оркатта: шаги 2--4 повторяются, пока
$\widehat\rho$ меняется.}
\end{seminarvisual}}

\newcommand{\runsvisual}{%
\begin{seminarvisual}
\begin{tikzpicture}[x=.57cm,y=.62cm,font=\small]
  % Problem 6.6: signs of OLS residuals e_1..e_24 and runs.
  \foreach \i/\sg in {1/m,2/m,3/m,4/m,5/m,6/m,7/p,8/m,9/p,10/p,11/p,12/p,13/m,14/p,15/m,16/p,17/p,18/m,19/m,20/p,21/p,22/p,23/p,24/p} {
    \if\sg p
      \fill[solbar!22] (\i-1,0) rectangle (\i,1); \node at (\i-.5,.5) {$+$};
    \else
      \fill[accent!18] (\i-1,0) rectangle (\i,1); \node at (\i-.5,.5) {$-$};
    \fi
    \node[text=gray,font=\scriptsize] at (\i-.5,1.35) {\i};
    \draw[white,thick] (\i-1,0)--(\i-1,1);
  }
  \foreach \a/\b/\k in {1/6/1,7/7/2,8/8/3,9/12/4,13/13/5,14/14/6,15/15/7,16/17/8,18/19/9,20/24/10} {
    \draw[notebar,thick] (\a-1+.1,-.12)--(\a-1+.1,-.3)--(\b-.1,-.3)--(\b-.1,-.12);
    \node[text=notebar,font=\scriptsize\bfseries] at ({(\a+\b-1)/2},-.68) {\k};
  }
  \node[anchor=east] at (-.15,.5) {знак $e_t$};
  \node[anchor=east,text=gray,font=\scriptsize] at (-.15,1.35) {$t$};
  \node[anchor=east,text=notebar] at (-.15,-.68) {ряд};
\end{tikzpicture}
\visualcaption{Задача 6.6: знаки остатков МНК в порядке времени. Скобки
отмечают ряды одинаковых знаков; их $k=10$. Плюсов $n_1=13$, минусов
$n_2=11$.}
\end{seminarvisual}}

\newcommand{\runsexamplevisual}{%
\begin{seminarvisual}
\begin{tikzpicture}[x=.5cm,y=.55cm,font=\small]
  % All 6 arrangements of n1 = 2 pluses and n2 = 2 minuses; brackets mark runs.
  \begin{scope}[shift={(0.00,0.00)}]
    \fill[solbar!22] (0,0) rectangle (1,1); \node at (0.5,.5) {$+$};
    \fill[solbar!22] (1,0) rectangle (2,1); \node at (1.5,.5) {$+$};
    \fill[accent!18] (2,0) rectangle (3,1); \node at (2.5,.5) {$-$};
    \fill[accent!18] (3,0) rectangle (4,1); \node at (3.5,.5) {$-$};
    \draw[white,thick] (1,0)--(1,1);
    \draw[white,thick] (2,0)--(2,1);
    \draw[white,thick] (3,0)--(3,1);
    \draw[notebar,thick] (0.10,-.12)--(0.10,-.32)--(1.90,-.32)--(1.90,-.12);
    \draw[notebar,thick] (2.10,-.12)--(2.10,-.32)--(3.90,-.32)--(3.90,-.12);
    \node[font=\small] at (2,-.85) {$k=2$};
  \end{scope}
  \begin{scope}[shift={(5.56,0.00)}]
    \fill[solbar!22] (0,0) rectangle (1,1); \node at (0.5,.5) {$+$};
    \fill[accent!18] (1,0) rectangle (2,1); \node at (1.5,.5) {$-$};
    \fill[solbar!22] (2,0) rectangle (3,1); \node at (2.5,.5) {$+$};
    \fill[accent!18] (3,0) rectangle (4,1); \node at (3.5,.5) {$-$};
    \draw[white,thick] (1,0)--(1,1);
    \draw[white,thick] (2,0)--(2,1);
    \draw[white,thick] (3,0)--(3,1);
    \draw[notebar,thick] (0.10,-.12)--(0.10,-.32)--(0.90,-.32)--(0.90,-.12);
    \draw[notebar,thick] (1.10,-.12)--(1.10,-.32)--(1.90,-.32)--(1.90,-.12);
    \draw[notebar,thick] (2.10,-.12)--(2.10,-.32)--(2.90,-.32)--(2.90,-.12);
    \draw[notebar,thick] (3.10,-.12)--(3.10,-.32)--(3.90,-.32)--(3.90,-.12);
    \node[font=\small] at (2,-.85) {$k=4$};
  \end{scope}
  \begin{scope}[shift={(11.12,0.00)}]
    \fill[solbar!22] (0,0) rectangle (1,1); \node at (0.5,.5) {$+$};
    \fill[accent!18] (1,0) rectangle (2,1); \node at (1.5,.5) {$-$};
    \fill[accent!18] (2,0) rectangle (3,1); \node at (2.5,.5) {$-$};
    \fill[solbar!22] (3,0) rectangle (4,1); \node at (3.5,.5) {$+$};
    \draw[white,thick] (1,0)--(1,1);
    \draw[white,thick] (2,0)--(2,1);
    \draw[white,thick] (3,0)--(3,1);
    \draw[notebar,thick] (0.10,-.12)--(0.10,-.32)--(0.90,-.32)--(0.90,-.12);
    \draw[notebar,thick] (1.10,-.12)--(1.10,-.32)--(2.90,-.32)--(2.90,-.12);
    \draw[notebar,thick] (3.10,-.12)--(3.10,-.32)--(3.90,-.32)--(3.90,-.12);
    \node[font=\small] at (2,-.85) {$k=3$};
  \end{scope}
  \begin{scope}[shift={(16.68,0.00)}]
    \fill[accent!18] (0,0) rectangle (1,1); \node at (0.5,.5) {$-$};
    \fill[solbar!22] (1,0) rectangle (2,1); \node at (1.5,.5) {$+$};
    \fill[solbar!22] (2,0) rectangle (3,1); \node at (2.5,.5) {$+$};
    \fill[accent!18] (3,0) rectangle (4,1); \node at (3.5,.5) {$-$};
    \draw[white,thick] (1,0)--(1,1);
    \draw[white,thick] (2,0)--(2,1);
    \draw[white,thick] (3,0)--(3,1);
    \draw[notebar,thick] (0.10,-.12)--(0.10,-.32)--(0.90,-.32)--(0.90,-.12);
    \draw[notebar,thick] (1.10,-.12)--(1.10,-.32)--(2.90,-.32)--(2.90,-.12);
    \draw[notebar,thick] (3.10,-.12)--(3.10,-.32)--(3.90,-.32)--(3.90,-.12);
    \node[font=\small] at (2,-.85) {$k=3$};
  \end{scope}
  \begin{scope}[shift={(22.24,0.00)}]
    \fill[accent!18] (0,0) rectangle (1,1); \node at (0.5,.5) {$-$};
    \fill[solbar!22] (1,0) rectangle (2,1); \node at (1.5,.5) {$+$};
    \fill[accent!18] (2,0) rectangle (3,1); \node at (2.5,.5) {$-$};
    \fill[solbar!22] (3,0) rectangle (4,1); \node at (3.5,.5) {$+$};
    \draw[white,thick] (1,0)--(1,1);
    \draw[white,thick] (2,0)--(2,1);
    \draw[white,thick] (3,0)--(3,1);
    \draw[notebar,thick] (0.10,-.12)--(0.10,-.32)--(0.90,-.32)--(0.90,-.12);
    \draw[notebar,thick] (1.10,-.12)--(1.10,-.32)--(1.90,-.32)--(1.90,-.12);
    \draw[notebar,thick] (2.10,-.12)--(2.10,-.32)--(2.90,-.32)--(2.90,-.12);
    \draw[notebar,thick] (3.10,-.12)--(3.10,-.32)--(3.90,-.32)--(3.90,-.12);
    \node[font=\small] at (2,-.85) {$k=4$};
  \end{scope}
  \begin{scope}[shift={(27.80,0.00)}]
    \fill[accent!18] (0,0) rectangle (1,1); \node at (0.5,.5) {$-$};
    \fill[accent!18] (1,0) rectangle (2,1); \node at (1.5,.5) {$-$};
    \fill[solbar!22] (2,0) rectangle (3,1); \node at (2.5,.5) {$+$};
    \fill[solbar!22] (3,0) rectangle (4,1); \node at (3.5,.5) {$+$};
    \draw[white,thick] (1,0)--(1,1);
    \draw[white,thick] (2,0)--(2,1);
    \draw[white,thick] (3,0)--(3,1);
    \draw[notebar,thick] (0.10,-.12)--(0.10,-.32)--(1.90,-.32)--(1.90,-.12);
    \draw[notebar,thick] (2.10,-.12)--(2.10,-.32)--(3.90,-.32)--(3.90,-.12);
    \node[font=\small] at (2,-.85) {$k=2$};
  \end{scope}
\end{tikzpicture}
\visualcaption{Все 6 расстановок двух плюсов и двух минусов; скобки отмечают
ряды. При $H_0$ каждая расстановка встречается с вероятностью $1/6$.}
\end{seminarvisual}}

\section{Основные формулы}

\textbf{Модель и предпосылка.} Рассматриваем регрессию
\[
  y_t = \beta_1 + \beta_2x_t + \varepsilon_t, \qquad t = 1, \dots, n,
\]
где $y_t$ — зависимая переменная, $x_t$ — регрессор, $\beta_1$, $\beta_2$ —
неизвестные параметры, $\varepsilon_t$ — ненаблюдаемая ошибка, $n$ — число
наблюдений, $t$ — время. Остатки МНК
$e_t = y_t - \widehat\beta_1 - \widehat\beta_2x_t$ служат оценками ошибок.
Одна из предпосылок МНК: ошибки разных наблюдений некоррелированы,
$\Cov(\varepsilon_t, \varepsilon_s) = 0$ при $t \ne s$. Её нарушение
называют \emph{автокорреляцией}; обычно она встречается во временных рядах.

\textbf{Модель AR(1).} Простейший вид автокорреляции — автокорреляция первого
порядка:
\[
  \varepsilon_t = \rho\varepsilon_{t-1} + \eta_t, \qquad |\rho| < 1,
\]
где $\rho$ — коэффициент автокорреляции (корреляция соседних ошибок), $\eta_t$
— независимые ошибки с нулевым средним и постоянной дисперсией (новая часть
ошибки в момент $t$). При $\rho > 0$ соседние ошибки обычно одного знака, при
$\rho < 0$ знаки чаще чередуются, при $\rho = 0$ автокорреляции нет.

\autocorrvisual

\textbf{Последствия.} Оценки МНК остаются несмещёнными, но перестают быть
эффективными, а обычные стандартные ошибки, $t$-тесты и доверительные
интервалы становятся неверными.

% -------------------------------------------------------------------
\subsection{Как обнаружить автокорреляцию}

Во всех тестах $H_0$: автокорреляции нет. Остатки берём строго в порядке
времени.

\textbf{1) График остатков} $e_t$ против $t$. Длинные участки одного знака —
признак положительной автокорреляции, частое чередование знаков —
отрицательной. Это подсказка, а не тест.

\textbf{2) Метод рядов.} Выписываем знаки остатков $e_1, \dots, e_n$ по
порядку. \emph{Ряд} — группа подряд идущих одинаковых знаков; например, в
последовательности $-\,-\,+\,+\,+\,-$ три ряда. Обозначим $k$ — число рядов,
$n_1$ — число плюсов, $n_2$ — число минусов. Если $H_0$ верна, знаки
расположены случайно, и
\begin{equation}
  M(k) = \frac{2n_1n_2}{n_1 + n_2} + 1,
  \qquad
  D(k) = \frac{2n_1n_2(2n_1n_2 - n_1 - n_2)}{(n_1 + n_2)^2(n_1 + n_2 - 1)},
  \label{eq:runs}
\end{equation}
где $M(k)$ и $D(k)$ — математическое ожидание и дисперсия числа рядов.
Наблюдаемое $k$ в эти формулы не входит: они зависят только от $n_1$ и $n_2$
и показывают, каким должно быть $k$, если $H_0$ верна. При $H_0$ все
расстановки $n_1$ плюсов и $n_2$ минусов по времени равновероятны, и $M(k)$,
$D(k)$ — среднее и дисперсия $k$ по всем таким расстановкам. Например, при
$n_1 = n_2 = 2$ возможны 6 расстановок:

\runsexamplevisual
\noindent Среднее и дисперсия $k$ по этим шести равновероятным расстановкам:
\begin{align*}
  M(k) &= \frac{2+4+3+3+4+2}{6} = 3,\\
  D(k) &= \frac{(2-3)^2+(4-3)^2+(3-3)^2+(3-3)^2+(4-3)^2+(2-3)^2}{6}
        = \frac46 = \frac23 .
\intertext{То же дают формулы \eqref{eq:runs}:}
  M(k) &= \frac{2n_1n_2}{n_1 + n_2} + 1 = \frac{2\cdot2\cdot2}{2+2} + 1 = 3,\\
  D(k) &= \frac{2n_1n_2(2n_1n_2 - n_1 - n_2)}{(n_1 + n_2)^2(n_1 + n_2 - 1)}
        = \frac{8\cdot(8 - 2 - 2)}{4^2\cdot 3} = \frac{32}{48} = \frac23 .
\end{align*}
Затем фактическое $k$ сравнивают с этим распределением. При $n_1, n_2 > 10$ число рядов
приближённо нормально, и $H_0$ не отвергается, если
\[
  k_1 = M(k) - z_{\alpha/2}\sqrt{D(k)} < k < M(k) + z_{\alpha/2}\sqrt{D(k)} = k_2 ,
\]
где $z_{\alpha/2}$ — критическое значение стандартного нормального
распределения ($z_{0{,}025} = 1{,}96$ при уровне 5\,\%). Слишком мало рядов
($k \le k_1$) — положительная автокорреляция, слишком много ($k \ge k_2$) —
отрицательная. При малых $n_1$, $n_2$ границы $k_1$, $k_2$ берут из таблицы
Сведа--Эйзенхарта. Тест использует только знаки остатков, но не их величину.

\textbf{3) Тест Дарбина--Уотсона.}
\begin{equation}
  \mathrm{DW} = \frac{\sum_{t=2}^{n}(e_t - e_{t-1})^2}{\sum_{t=1}^{n}e_t^2}
  \approx 2(1 - \widehat{\rho}),
  \qquad
  \widehat{\rho} = \frac{\sum_{t=2}^{n}e_te_{t-1}}{\sum_{t=2}^{n}e_{t-1}^2},
  \label{eq:dw}
\end{equation}
где $\widehat\rho$ — оценка $\rho$ по остаткам (МНК-оценка регрессии $e_t$ на
$e_{t-1}$ без константы).

\noindent\emph{Откуда $\mathrm{DW} \approx 2(1 - \widehat\rho)$.} Раскроем
квадрат в числителе:
\[
  \sum_{t=2}^{n}(e_t - e_{t-1})^2
  = \sum_{t=2}^{n}e_t^2 + \sum_{t=2}^{n}e_{t-1}^2 - 2\sum_{t=2}^{n}e_te_{t-1} .
\]
Каждое из трёх слагаемых приближённо выражается через $\sum_{t=1}^{n}e_t^2$
(знаменатель DW):
\begin{itemize}[leftmargin=2.2em, itemsep=1pt]
  \item $\sum_{t=2}^{n}e_t^2 = e_2^2 + \dots + e_n^2$ — это вся сумма
        $\sum_{t=1}^{n}e_t^2$ без одного слагаемого $e_1^2$, поэтому
        $\approx \sum_{t=1}^{n}e_t^2$;
  \item $\sum_{t=2}^{n}e_{t-1}^2 = e_1^2 + \dots + e_{n-1}^2$ — вся сумма без
        $e_n^2$, поэтому тоже $\approx \sum_{t=1}^{n}e_t^2$;
  \item из определения $\widehat\rho$ в \eqref{eq:dw}
        $\sum_{t=2}^{n}e_te_{t-1} = \widehat\rho\sum_{t=2}^{n}e_{t-1}^2
        \approx \widehat\rho\sum_{t=1}^{n}e_t^2$.
\end{itemize}
Одно лишнее слагаемое мало по сравнению с суммой из $n$ слагаемых, поэтому
при большом $n$ эти приближения точные. Подставляем:
\[
  \mathrm{DW} \approx \frac{\sum e_t^2 + \sum e_t^2 - 2\widehat\rho\sum e_t^2}{\sum e_t^2}
  = \frac{2\sum e_t^2 - 2\widehat\rho\sum e_t^2}{\sum e_t^2} = 2(1 - \widehat\rho).
\]
Поэтому $\mathrm{DW} \approx 0$ при $\widehat\rho \approx 1$,
$\mathrm{DW} \approx 2$ при $\widehat\rho \approx 0$ и $\mathrm{DW} \approx 4$
при $\widehat\rho \approx -1$. Решение принимают по табличным границам
$d_L < d_U$, которые зависят от $n$ и числа регрессоров без константы:
\begin{center}
\begin{tabular*}{\linewidth}{@{\extracolsep{\fill}}ll@{}}
\toprule
Значение DW & Вывод \\
\midrule
$0 \le \mathrm{DW} < d_L$ & положительная автокорреляция \\
$d_L \le \mathrm{DW} \le d_U$ & вывода нет \\
$d_U < \mathrm{DW} < 4 - d_U$ & автокорреляции нет \\
$4 - d_U \le \mathrm{DW} \le 4 - d_L$ & вывода нет \\
$4 - d_L < \mathrm{DW} \le 4$ & отрицательная автокорреляция \\
\bottomrule
\end{tabular*}
\end{center}
\noindent На рисунке~\ref{fig:dwzones} те же пять зон показаны на отрезке
$[0;\,4]$, на котором лежат все возможные значения DW: у 0 —
положительная автокорреляция, около 2 — её нет, у 4 — отрицательная, а
между ними — две зоны, где вывода сделать нельзя.

\dwzonesvisual

Таблицы DW применимы к регрессии с константой, в которой среди регрессоров
нет лагированной зависимой переменной $y_{t-1}$.

% -------------------------------------------------------------------
\subsection{Коррекция: процедура Кохрейна--Оркатта}

Если $\rho$ известно, вычтем из уравнения для момента $t$ то же уравнение
для момента $t-1$, умноженное на $\rho$:
\begin{equation}
  y_t - \rho y_{t-1} = \beta_1(1 - \rho) + \beta_2(x_t - \rho x_{t-1}) + \eta_t,
  \qquad t = 2, \dots, n .
  \label{eq:co}
\end{equation}
Величины $y_t - \rho y_{t-1}$ и $x_t - \rho x_{t-1}$ называют
\emph{квазиразностями}. Ошибка $\eta_t$ уже не автокоррелирована, поэтому
уравнение оценивают обычным МНК с константой; первое наблюдение теряется.
Константа этого уравнения равна $\beta_1(1 - \rho)$, поэтому
$\widehat\beta_1$ получают делением её оценки на $1 - \rho$, а наклон сразу
даёт $\widehat\beta_2$. На практике $\rho$ неизвестно, и его оценивают по
остаткам, уточняя по кругу:

\coloopvisual
\fi

% ===================================================================
\ifsolutions\clearpage\fi
\section{Задачи}

% -------------------------------------------------------------------
\problem{6.6}

Рассмотрим модель, связывающую количество вакансий $w_t$ и уровень
безработицы $u_t$:
\[
  \ln w_t = \beta_1 + \beta_2\ln u_t + \varepsilon_t .
\]
Ошибки $\varepsilon_t$ независимы и нормально распределены
$N(0, \sigma_{\varepsilon}^2)$.
\begin{enumerate}[label=\asbuk*), leftmargin=2.2em, itemsep=3pt]
  \item Используя (искусственные) данные из таблицы, найдите МНК-оценки
        параметров $\beta_1$ и $\beta_2$, а также 95\,\%-ный доверительный
        интервал для $\beta_2$.
  \item Вычислите статистику Дарбина--Уотсона и проверьте отсутствие
        автокорреляции также методом рядов. Что это говорит об
        исходном предположении об ошибках $\varepsilon_t$? Что можно сказать
        о доверительном интервале, найденном в п.~а)?
  \item Оцените модель заново, используя модель автокорреляции первого
        порядка для ошибок регрессии. Найдите 95\,\%-ный доверительный
        интервал для $\beta_2$. Сравните результат с интервалом, полученным в
        п.~а).
\end{enumerate}

\begin{center}
\small
\begin{tabular*}{\linewidth}{@{\extracolsep{\fill}}rccrccrcc@{}}
\toprule
$t$ & $w_t$ & $u_t$ & $t$ & $w_t$ & $u_t$ & $t$ & $w_t$ & $u_t$ \\
\midrule
1 & 1,73 & 8,65 &  9 & 5,06 & 2,87 & 17 & 3,15 & 4,72 \\
2 & 1,94 & 4,82 & 10 & 2,81 & 5,29 & 18 & 1,92 & 7,45 \\
3 & 3,05 & 2,67 & 11 & 4,43 & 3,31 & 19 & 2,26 & 6,21 \\
4 & 4,17 & 2,67 & 12 & 3,19 & 5,44 & 20 & 6,18 & 2,64 \\
5 & 2,52 & 2,58 & 13 & 2,23 & 6,80 & 21 & 2,07 & 8,55 \\
6 & 1,71 & 8,07 & 14 & 2,06 & 8,25 & 22 & 8,39 & 2,60 \\
7 & 1,95 & 8,83 & 15 & 3,33 & 3,44 & 23 & 2,75 & 6,25 \\
8 & 2,57 & 5,54 & 16 & 2,12 & 7,80 & 24 & 6,10 & 2,70 \\
\bottomrule
\end{tabular*}
\end{center}

\begin{excel}
Добавьте столбцы \xl{=LN(w)} и \xl{=LN(u)} и оцените регрессию
\ui{Анализ данных} $\to$ \ui{Регрессия} (галочка \ui{Остатки}). DW по
остаткам $e_1, \dots, e_{24}$: \xl{=СУММКВРАЗН($e_2{:}e_{24}$; $e_1{:}e_{23}$)/СУММКВ($e_1{:}e_{24}$)}.
Метод рядов: столбец знаков \xl{=ЕСЛИ($e_t$>0;1;0)}, число рядов
\xl{=1+СУММПРОИЗВ(--($s_2{:}s_{24}$<>$s_1{:}s_{23}$))}, где $s_t$ — знаки.
Для п.~в) — столбцы $y_t - \widehat{\rho}y_{t-1}$ и
$x_t - \widehat{\rho}x_{t-1}$ и регрессия первого на второй с константой. Это один шаг коррекции; для следующего шага пересчитайте остатки исходного уравнения.
\end{excel}

\begin{solution}
Обозначим $y_t=\ln w_t$, $x_t=\ln u_t$; $n=24$. Наблюдения не сортируем:
DW и метод рядов используют порядок времени. План: МНК и интервал (п.~а),
проверка независимости ошибок (п.~б), пересчёт с учётом AR(1) (п.~в).

\medskip
\textbf{а) МНК.}
\begin{center}
\begin{tabular*}{\linewidth}{@{\extracolsep{\fill}}lrrrr@{}}
\toprule
Переменная & Коэффициент & Ст. ошибка & $t$-статистика & $P$-значение \\
\midrule
const   & $\phantom{-}2{,}2997$ & 0,1860 & $12{,}366$ & $<0{,}0001$ \\
$\ln u$ & $-0{,}7791$ & 0,1133 & $-6{,}875$ & $<0{,}0001$ \\
\midrule
\multicolumn{5}{l}{$R^2 = 0{,}6824$;\quad $\ESS = 1{,}4119$;\quad $\mathrm{DW} = 1{,}1026$} \\
\bottomrule
\end{tabular*}
\end{center}
$\widehat{\beta}_1 = 2{,}30$, $\widehat{\beta}_2 = -0{,}78$: рост безработицы
на 1\,\% сопровождается в среднем снижением числа вакансий на 0,78\,\%.
Доверительный интервал по обычной формуле:
\[
  \widehat{\beta}_2 \pm t_{0{,}025}(n-2)\,s_{\widehat{\beta}_2}
  = -0{,}7791 \pm 2{,}074\cdot 0{,}1133 = (-1{,}014;\ -0{,}544),
\]
где $s_{\widehat\beta_2} = 0{,}1133$ — стандартная ошибка наклона из
таблицы, $t_{0{,}025}(22) = 2{,}074$ — критическое значение при
$n - 2 = 22$ степенях свободы.

\medskip
\textbf{б) Статистика Дарбина--Уотсона.} По остаткам $e_t$ \eqref{eq:dw}:
\[
  \mathrm{DW} = \frac{\sum_{t=2}^{24}(e_t - e_{t-1})^2}{\sum_{t=1}^{24}e_t^2}
              = \frac{1{,}5568}{1{,}4119} = 1{,}103,
  \qquad
  \widehat{\rho} \approx 1 - \frac{\mathrm{DW}}{2} = 0{,}45 ,
\]
где оценка $\widehat\rho$ получена из приближения
$\mathrm{DW} \approx 2(1 - \widehat\rho)$ в \eqref{eq:dw}. По таблице DW
для $n = 24$ и одного регрессора на уровне 5\,\% $d_L = 1{,}27$,
$d_U = 1{,}45$. $\mathrm{DW} = 1{,}10 < d_L$, поэтому гипотеза об отсутствии
автокорреляции отвергается: есть положительная автокорреляция первого
порядка. Значит, исходное предположение о \emph{независимости} ошибок
нарушено. Оценки МНК остаются несмещёнными, но стандартные ошибки из п.~а)
посчитаны по формулам, верным только для независимых ошибок, и поэтому
ненадёжны. Следовательно, интервал из п.~а) тоже ненадёжен: его настоящий
уровень доверия может заметно отличаться от 95\,\%.

\dwvisual

\smallskip
\noindent\begin{minipage}{\linewidth}
\parindent=1.5em
\emph{Проверка методом рядов.} Выписываем знаки остатков $e_t$ из п.~а) в
порядке времени и отмечаем ряды:

\runsvisual
\end{minipage}
Рядов $k = 10$, плюсов $n_1 = 13$, минусов $n_2 = 11$. По \eqref{eq:runs}
\[
  M(k) = \frac{2n_1n_2}{n_1 + n_2} + 1 = \frac{2\cdot 13\cdot 11}{24} + 1 = 12{,}92,
\]
\[
  D(k) = \frac{2n_1n_2(2n_1n_2 - n_1 - n_2)}{(n_1 + n_2)^2(n_1 + n_2 - 1)}
       = \frac{286\cdot(286 - 24)}{24^2\cdot 23} = \frac{74\,932}{13\,248} = 5{,}656,
\]
откуда $\sqrt{D(k)} = 2{,}378$.
Границы на уровне 5\,\%:
\[
  k_1 = 12{,}92 - 1{,}96\cdot 2{,}378 = 8{,}26,
  \qquad
  k_2 = 12{,}92 + 1{,}96\cdot 2{,}378 = 17{,}58 .
\]
$8{,}26 < k = 10 < 17{,}58$, поэтому метод рядов гипотезу об отсутствии
автокорреляции \emph{не отвергает}. Таблица Сведа--Эйзенхарта для
$n_1 = 13$, $n_2 = 11$ даёт нижнюю границу $k_1 = 7$ — вывод тот же.

Тесты расходятся, потому что метод рядов смотрит только на знаки остатков.
Остатки $e_8 = -0{,}022$ и $e_{13} = -0{,}004$ почти нулевые, и их знак
почти случаен; каждый из них разрывает длинный ряд плюсов и добавляет два
ряда. Будь они положительными, рядов было бы 6 — меньше нижней границы
(около 7,9 при $n_1 = 15$, $n_2 = 9$), и $H_0$ отверглась бы. Статистика DW
учитывает и величину остатков: в начале выборки идут крупные отрицательные
остатки, в конце — крупные положительные, поэтому разности $e_t - e_{t-1}$
малы. Здесь DW — более чувствительный тест, и вывод п.~б) остаётся прежним:
автокорреляция есть.

\medskip
\textbf{в) Оценка с учётом AR(1).} В п.~б) выяснилось, что соседние ошибки
связаны. Опишем эту связь моделью AR(1):
\[
  \varepsilon_t = \rho\varepsilon_{t-1} + \eta_t ,
\]
где $\rho$ — коэффициент автокорреляции ($|\rho| < 1$), $\eta_t$ —
независимые ошибки с постоянной дисперсией. Идея: преобразовать уравнение
так, чтобы в нём осталась только ошибка $\eta_t$. Тогда предпосылки МНК
выполнены, и обычные стандартные ошибки и интервал снова верны.

\smallskip
\emph{Шаг 1. Преобразование (пока считаем $\rho$ известным).} Запишем модель
для моментов $t$ и $t-1$, второе уравнение умножим на $\rho$:
\[
  y_t = \beta_1 + \beta_2x_t + \varepsilon_t ,
  \qquad
  \rho y_{t-1} = \rho\beta_1 + \rho\beta_2x_{t-1} + \rho\varepsilon_{t-1} .
\]
Вычтем второе из первого. Ошибка становится
$\varepsilon_t - \rho\varepsilon_{t-1} = \eta_t$, и получаем модель в
квазиразностях \eqref{eq:co}:
\[
  {\underbrace{y_t - \rho y_{t-1}}_{y_t^*}}
  = {\underbrace{\beta_1(1 - \rho)}_{\beta_1^*}}
  + \beta_2{\underbrace{(x_t - \rho x_{t-1})}_{x_t^*}} + \eta_t,
  \qquad t = 2, \dots, 24 ,
\]
где $y_t^*$, $x_t^*$ — новые переменные, $\beta_1^*$ — свободный член
нового уравнения. Это обычная регрессия $y_t^*$ на $x_t^*$ с константой.
Её наклон — это сразу $\beta_2$, а $\beta_1$ находится из константы:
$\beta_1 = \beta_1^*/(1 - \rho)$. Наблюдение $t = 1$ теряется: для него нет
$y_0$ и $x_0$, поэтому остаётся 23 наблюдения.

\smallskip
\emph{Шаг 2. Оценка $\rho$.} На деле $\rho$ неизвестно. Уравнение
$\varepsilon_t = \rho\varepsilon_{t-1} + \eta_t$ — регрессия $\varepsilon_t$
на $\varepsilon_{t-1}$ без константы. Заменяем ошибки остатками $e_t$ из
п.~а) и получаем МНК-оценку \eqref{eq:dw}:
\[
  \widehat\rho = \frac{\sum_{t=2}^{24} e_te_{t-1}}{\sum_{t=2}^{24} e_{t-1}^2}
  = \frac{0{,}5911}{1{,}3321} = 0{,}4437 .
\]
Это близко к грубой оценке $1 - \mathrm{DW}/2 = 0{,}45$ из п.~б).

\smallskip
\emph{Шаг 3. МНК для квазиразностей с $\widehat\rho = 0{,}4437$.} Для
каждого $t = 2, \dots, 24$ считаем новые переменные, например для $t = 2$:
\begin{align*}
  y_2^* &= y_2 - 0{,}4437\,y_1 = 0{,}6627 - 0{,}4437\cdot 0{,}5481 = 0{,}4195 ,\\
  x_2^* &= x_2 - 0{,}4437\,x_1 = 1{,}5728 - 0{,}4437\cdot 2{,}1576 = 0{,}6154 .
\end{align*}
Регрессия $y_t^*$ на $x_t^*$ по 23 наблюдениям даёт
$\widehat\beta_1^* = 1{,}3302$ и $\widehat\beta_2 = -0{,}8296$, откуда
$\widehat\beta_1 = 1{,}3302/(1 - 0{,}4437) = 2{,}3913$.

\smallskip
\emph{Шаг 4. Итерации Кохрейна--Оркатта.} Оценка $\widehat\rho$ из шага~2
построена по остаткам обычного МНК, который автокорреляцию не учитывал.
Теперь у нас есть уточнённые $\widehat\beta_1$, $\widehat\beta_2$. По ним
считаем новые остатки \emph{исходного} уравнения
$e_t = y_t - \widehat\beta_1 - \widehat\beta_2x_t$, заново оцениваем
$\widehat\rho$ (шаг~2) и повторяем шаг~3. Останавливаемся, когда
$\widehat\rho$ перестаёт меняться:
\begin{center}
\begin{tabular*}{\linewidth}{@{\extracolsep{\fill}}ccccc@{}}
\toprule
Итерация & $\widehat\rho$ & $\widehat\beta_1^*$ & $\widehat\beta_1 = \widehat\beta_1^*/(1-\widehat\rho)$ & $\widehat\beta_2$ \\
\midrule
1 & 0,4437 & 1,3302 & 2,3913 & $-0{,}8296$ \\
2 & 0,4668 & 1,2768 & 2,3948 & $-0{,}8312$ \\
3 & 0,4675 & 1,2752 & 2,3949 & $-0{,}8312$ \\
4 & 0,4676 & 1,2751 & 2,3949 & $-0{,}8312$ \\
\bottomrule
\end{tabular*}
\end{center}

\smallskip
\emph{Шаг 5. Итоговые оценки.} Последняя регрессия:
$y_t - 0{,}4676y_{t-1}$ на $x_t - 0{,}4676x_{t-1}$. Её константа
$\widehat\beta_1^* = 1{,}2751$ со стандартной ошибкой $0{,}0889$. Оба числа
делим на $1 - \widehat\rho = 0{,}5324$, чтобы вернуться к $\beta_1$:
$\widehat\beta_1 = 2{,}3949$, $s_{\widehat\beta_1} = 0{,}0889/0{,}5324 =
0{,}1669$. Наклон и его ошибка пересчёта не требуют. Стандартные ошибки
посчитаны так, будто $\widehat\rho$ известно точно:
\begin{center}
\begin{tabular*}{\linewidth}{@{\extracolsep{\fill}}lrrrr@{}}
\toprule
Переменная & Коэффициент & Ст. ошибка & $t$-статистика & $P$-значение \\
\midrule
const   & $\phantom{-}2{,}3949$ & 0,1669 & $14{,}349$ & $<0{,}0001$ \\
$\ln u$ & $-0{,}8312$ & 0,0930 & $-8{,}939$ & $<0{,}0001$ \\
\midrule
\multicolumn{5}{l}{$\widehat\rho=0{,}4676$;\quad $\mathrm{DW} = 2{,}1617$} \\
\bottomrule
\end{tabular*}
\end{center}

\noindent\begin{minipage}{\linewidth}
\emph{Шаг 6. Интервал.} Приближённый интервал для $\beta_2$ по конечной регрессии
(23 наблюдения, константа и наклон: $23-2=21$ степень свободы,
$t_{0{,}025}(21)=2{,}080$):
\[
  -0{,}8312 \pm 2{,}080\cdot 0{,}0930 = (-1{,}025;\ -0{,}638).
\]
Это приближение: дополнительная неопределённость оценки $\rho$
в обычной формуле стандартной ошибки последнего шага не учитывается.
\end{minipage}

\medskip
\noindent\begin{minipage}{\linewidth}
\textbf{Сравнение.}
\intervalvisual
Оценка наклона изменилась мало ($-0{,}83$ против $-0{,}78$): рост
безработицы на 1\,\% связан со снижением числа вакансий примерно на
0,83\,\%. Интервал с учётом AR(1) уже (длина 0,39 против 0,47). Но это не
доказывает, что точность выросла: интервал из п.~а) построен при нарушенной
предпосылке, и его ширине нельзя доверять. Второй интервал верен
приближённо и только если ошибки действительно следуют AR(1). Проверка:
у остатков последней регрессии $\widehat\eta_t$ статистика
$\mathrm{DW} = 2{,}16$, близко к 2, — автокорреляции первого порядка в них
не видно (графики ниже).
\end{minipage}
\arresidualvisual
\end{solution}

% -------------------------------------------------------------------
\ifsolutions\clearpage\fi
\problem{6.3}

Рассмотрим уравнение $y_t = \alpha + \beta x_t + \varepsilon_t$, где ошибки
$\varepsilon_t$ порождаются авторегрессионным процессом второго порядка:
\[
  \varepsilon_t = \rho_1\varepsilon_{t-1} + \rho_2\varepsilon_{t-2} + u_t .
\]
Предложите обобщение итеративной процедуры Кохрейна--Оркатта для
оценивания параметров этой модели.

\begin{solution}
\textbf{Идея.} Здесь $\rho_1$, $\rho_2$ — коэффициенты автокорреляции,
$u_t$ — некоррелированные ошибки с постоянной дисперсией. Ошибка
$\varepsilon_t$ зависит уже от \emph{двух} предыдущих ошибок. Как и для
AR(1) в \eqref{eq:co}, преобразуем уравнение так, чтобы в нём осталась
только ошибка $u_t$; для этого придётся вычесть два лаговых уравнения, а не
одно.

\smallskip
\emph{Шаг 1. Преобразование (пока считаем $\rho_1$, $\rho_2$ известными).}
Запишем модель для моментов $t$, $t-1$ и $t-2$ и умножим второе уравнение на
$\rho_1$, третье — на $\rho_2$:
\begin{align*}
  y_t &= \alpha + \beta x_t + \varepsilon_t ,\\
  \rho_1y_{t-1} &= \rho_1\alpha + \rho_1\beta x_{t-1} + \rho_1\varepsilon_{t-1} ,\\
  \rho_2y_{t-2} &= \rho_2\alpha + \rho_2\beta x_{t-2} + \rho_2\varepsilon_{t-2} .
\end{align*}
Вычтем второе и третье из первого. Ошибка становится
$\varepsilon_t - \rho_1\varepsilon_{t-1} - \rho_2\varepsilon_{t-2} = u_t$, и
получаем
\[
  {\underbrace{y_t - \rho_1y_{t-1} - \rho_2y_{t-2}}_{y_t^*}}
  = {\underbrace{\alpha(1 - \rho_1 - \rho_2)}_{\alpha^*}}
  + \beta{\underbrace{(x_t - \rho_1x_{t-1} - \rho_2x_{t-2})}_{x_t^*}} + u_t,
  \qquad t = 3, \dots, n ,
\]
где $y_t^*$, $x_t^*$ — новые переменные, $\alpha^*$ — свободный член нового
уравнения. Ошибка $u_t$ некоррелирована, поэтому это обычная регрессия
$y_t^*$ на $x_t^*$ с константой. Её наклон — сразу $\beta$, а $\alpha$
находится из константы: $\alpha = \alpha^*/(1 - \rho_1 - \rho_2)$. Теряются
\emph{два} первых наблюдения: для $t = 1$ и $t = 2$ нет нужных лагов, остаётся
$n - 2$ наблюдения.

\smallskip
\emph{Шаг 2. МНК исходной модели.} На практике $\rho_1$, $\rho_2$ неизвестны.
Оцениваем $y_t = \alpha + \beta x_t + \varepsilon_t$ обычным МНК и находим
остатки $e_t = y_t - \widehat\alpha - \widehat\beta x_t$, $t = 1, \dots, n$.

\smallskip
\emph{Шаг 3. Оценка $\rho_1$, $\rho_2$.} Уравнение для ошибок — регрессия
$\varepsilon_t$ на $\varepsilon_{t-1}$ и $\varepsilon_{t-2}$ без константы.
Заменяем ошибки остатками и оцениваем МНК
\[
  e_t = \rho_1e_{t-1} + \rho_2e_{t-2} + v_t, \qquad t = 3, \dots, n ,
\]
где $v_t$ — ошибка этой вспомогательной регрессии. Это множественная
регрессия с двумя регрессорами $e_{t-1}$ и $e_{t-2}$; константы в ней нет,
потому что её нет и в модели ошибок. Получаем $\widehat\rho_1$,
$\widehat\rho_2$.

\smallskip
\emph{Шаг 4. МНК для преобразованного уравнения.} Подставляем
$\widehat\rho_1$, $\widehat\rho_2$ в шаг 1: считаем
$y_t^* = y_t - \widehat\rho_1y_{t-1} - \widehat\rho_2y_{t-2}$ и
$x_t^* = x_t - \widehat\rho_1x_{t-1} - \widehat\rho_2x_{t-2}$ для
$t = 3, \dots, n$, оцениваем регрессию $y_t^*$ на $x_t^*$ с константой и
получаем
\[
  \widehat\beta = \text{наклон},
  \qquad
  \widehat\alpha = \frac{\widehat\alpha^*}{1 - \widehat\rho_1 - \widehat\rho_2} .
\]

\smallskip
\emph{Шаг 5. Итерации.} По новым $\widehat\alpha$, $\widehat\beta$ считаем
остатки \emph{исходного} уравнения $e_t = y_t - \widehat\alpha -
\widehat\beta x_t$ для всех $t = 1, \dots, n$ и возвращаемся к шагу~3.
Останавливаемся, когда $\widehat\rho_1$ и $\widehat\rho_2$ перестают
меняться (разность соседних оценок меньше заданной точности).

\medskip
\noindent\begin{minipage}{\linewidth}
\textbf{Что меняется по сравнению с AR(1).}
\begin{center}
\small
\begin{tabular*}{\linewidth}{@{\extracolsep{\fill}}lll@{}}
\toprule
 & AR(1) & AR(2) \\
\midrule
Модель ошибок & $\varepsilon_t = \rho\varepsilon_{t-1} + \eta_t$
              & $\varepsilon_t = \rho_1\varepsilon_{t-1} + \rho_2\varepsilon_{t-2} + u_t$ \\
Оценка $\rho$ & $e_t$ на $e_{t-1}$ & $e_t$ на $e_{t-1}$, $e_{t-2}$ \\
Квазиразность & $y_t - \rho y_{t-1}$ & $y_t - \rho_1y_{t-1} - \rho_2y_{t-2}$ \\
Наблюдений в регрессии & $n - 1$ & $n - 2$ \\
Свободный член & $\alpha(1 - \rho)$ & $\alpha(1 - \rho_1 - \rho_2)$ \\
\bottomrule
\end{tabular*}
\end{center}
\end{minipage}
В обоих случаях наклон преобразованного уравнения сразу равен наклону
исходной модели, а свободный член надо разделить на $1 - \rho$ (или на
$1 - \rho_1 - \rho_2$).
\end{solution}
